\chapter{Satellieten en de beweging van de planeten}\label{ch:g12:satellites-kepler}

De schotel op het balkon zit vastgeschroefd en staart toch de hele dag naar
één vast punt van de hemel, \qty{36000}{km} hoog; het ruimtestation steekt
de avondhemel in minuten over; een zwerm vliegende klokken vertelt je
telefoon waar je staat. Geen van die machines heeft een motor draaien: ze
\emph{vallen} allemaal gewoon. Dit hoofdstuk rekent uit hoe snel en hoe hoog
elk daarvan moet vallen --- en weegt met dezelfde wet planeten.

\section{Vallen rond de Aarde}

Alles rust op de wet die we eerder hebben opgemeten
(\cref{ch:g10:universal-gravitation}): de Aarde, met massa $M$, trekt aan een
\omterm{def:g10:universal-gravitation:satellite}{satelliet} met massa $m$ waarvan het middelpunt op een afstand $r$ van het
hare ligt, met een \omterm{def:g10:inertia:force}{kracht} naar het middelpunt van de Aarde en met grootte
\[ F = G\,\frac{m M}{r^2}, \qquad G = \qty{6.67e-11}{N.m^2/kg^2}. \]
Let op wat $r$ is: de afstand \emph{tussen de middelpunten}. Een \omterm{def:g10:universal-gravitation:satellite}{satelliet} op
hoogte $h$ heeft $r = R + h$, met $R = \qty{6.37e6}{m}$ de straal van de
Aarde --- de $R$ vergeten is de klassieke blunder van dit hoofdstuk.

\begin{remark}[Het kanon van Newton, opnieuw]\label{rem:g12:satellites-kepler:cannon}
Horizontaal afgeschoten landt een kanonskogel; sneller landt hij verder; bij
de juiste snelheid kromt de grond precies even snel weg als de kogel valt, en
sluit de val zich: een \omterm{def:g10:universal-gravitation:satellite}{omloopbaan}. Een \omterm{def:g10:universal-gravitation:satellite}{satelliet} heeft geen motor nodig ---
alleen de juiste snelheid \emph{opzij}.
\end{remark}

\begin{omfigure}
\begin{tikzpicture}[scale=1.0, font=\small]
  \fill[omExo!15] (0,0) circle (1.5); \draw[omExo, thick] (0,0) circle (1.5);
  \fill[omExo!60] (-0.22,1.44) -- (0,1.85) -- (0.22,1.44) -- cycle;
  \coordinate (L) at (0,1.85);
  \draw[omProp, thick] (L) .. controls (0.7,1.85) .. (50:1.5);
  \draw[omProp, thick] (L) .. controls (1.4,1.85) and (2.0,0.9) .. (0:1.5);
  \draw[-Stealth, omDef, thick] (0,1.85) arc (90:-260:1.85);
  \draw[-Stealth, omThm, thick] (L) .. controls (1.7,2.0) and (3.2,2.2)
    .. (4.3,2.9);
  \node[omDef, anchor=east] at (-1.85,1.1) {omloopbaan};
  \node[omThm, anchor=west] at (3.35,2.35) {ontsnapt};
  \node[omProp, anchor=west] at (1.35,0.55) {valt};
\end{tikzpicture}

{\small Het kanon van Newton: steeds sneller afgeschoten landt de kogel
verder (oranje); bij ongeveer \qty{7.9}{km/s} volgt zijn val de kromming van
de Aarde en komt hij in een \omterm{def:g10:universal-gravitation:satellite}{omloopbaan} (blauw); nog sneller ontsnapt hij
(rood).}
\end{omfigure}

\section{De cirkelvormige omloopbaan: één snelheid, één omlooptijd}

\begin{theorem}[Baansnelheid]\label{thm:g12:satellites-kepler:speed}
Een \omterm{def:g10:universal-gravitation:satellite}{satelliet} in een eenparige cirkelvormige \omterm{def:g10:universal-gravitation:satellite}{omloopbaan} met straal $r$ om een
lichaam met massa $M$ beweegt met de snelheid
\[ v = \sqrt{\frac{G M}{r}}, \]
die niet van de massa van de \omterm{def:g10:universal-gravitation:satellite}{satelliet} afhangt.
\end{theorem}

\begin{proof}
De \omterm{def:g11:fundamental-interactions:four}{gravitatie} is de enige \omterm{def:g10:inertia:force}{kracht}, dus luidt de tweede wet van Newton
(\cref{thm:g12:newtons-laws:second}) $m\vect a = \vect F$; projecteer die op
het \omterm{def:g12:kinematics-2d:frenet}{frenetstelsel} (\cref{def:g12:kinematics-2d:frenet}). De \omterm{def:g10:inertia:force}{kracht} mikt op
het middelpunt: haar tangentiële component is nul, dus is de snelheid
constant; haar normale component is de volle $GmM/r^2$, en bij een
cirkelbeweging is $a_n = v^2/r$
(\cref{thm:g12:kinematics-2d:centripetal}). Dus $m v^2/r = GmM/r^2$, en ---
$m$ valt weg --- $v^2 = GM/r$.
\end{proof}

\begin{omfigure}
\begin{tikzpicture}[scale=1.0, font=\small]
  \fill[omExo!15] (0,0) circle (0.85); \draw[omExo, thick] (0,0) circle (0.85);
  \node at (0,0) {Aarde}; \draw[dashed, omExo] (0,0) circle (2.3);
  \draw[thin] (0,0) -- (40:2.3) node[midway, above left] {$r$};
  \fill (40:2.3) circle (2.2pt) node[above right] {$S$};
  \draw[-Stealth, omThm, very thick] (40:2.3) -- (40:1.35)
    node[pos=0.85, below right] {$\vect F$};
  \draw[-Stealth, omDef, thick] (40:2.3) -- ++(130:1.15)
    node[above] {$\vect v$};
\end{tikzpicture}

{\small Het hele mechanisme van een \omterm{def:g10:universal-gravitation:satellite}{omloopbaan}: de snelheid raakt de \omterm{def:g10:relative-motion:trajectory}{baan},
de \omterm{def:g10:inertia:force}{kracht} wijst naar het middelpunt --- de \omterm{def:g11:fundamental-interactions:four}{gravitatie} remt de \omterm{def:g10:universal-gravitation:satellite}{satelliet}
nooit af, ze buigt alleen zijn weg, voor altijd.}
\end{omfigure}

\begin{proposition}[Omlooptijd van een satelliet]\label{prop:g12:satellites-kepler:period}
De \omterm{def:g10:universal-gravitation:satellite}{satelliet} doorloopt zijn \omterm{def:g10:universal-gravitation:satellite}{omloopbaan} in de \omterm{def:g12:kinematics-2d:ucm}{omlooptijd}
\[ T = \frac{2\pi r}{v} = 2\pi \sqrt{\frac{r^3}{G M}}. \]
\end{proposition}

\begin{proof}
Eén ronde is de omtrek $2\pi r$ bij constante snelheid; vul
$v = \sqrt{GM/r}$ in: $T = 2\pi r\sqrt{r/(GM)} = 2\pi\sqrt{r^3/(GM)}$.
\end{proof}

Noch $v$ noch $T$ bevat $m$: een schroefje dat van het ruimtestation
loskomt, draait er gewoon naast mee. En beide dalen met $r$: \emph{hoe hoger
de \omterm{def:g10:universal-gravitation:satellite}{omloopbaan}, hoe trager de \omterm{def:g10:universal-gravitation:satellite}{satelliet}}, in snelheid
($v \propto 1/\sqrt r$) en in hoek ($T \propto r^{3/2}$).

\begin{example}[Het ruimtestation]\label{ex:g12:satellites-kepler:iss}
Het internationale ruimtestation vliegt op hoogte
$h \approx \qty{400}{km}$, dus $r = \num{6.37e6} + \num{4.0e5} =
\qty{6.77e6}{m}$; met $GM = \num{6.67e-11} \times \num{5.97e24} =
\qty{3.98e14}{m^3/s^2}$ is $v = \sqrt{\num{3.98e14}/\num{6.77e6}}
\approx \qty{7.7e3}{m/s}$ en $T = 2\pi r/v \approx \qty{5.5e3}{s}
\approx \qty{92}{min}$: elke anderhalf uur een ronde om de planeet, en zo'n
zestien zonsopgangen per dag voor de astronauten.
\end{example}

\begin{example}[De Maan valt zoals de appel]\label{ex:g12:satellites-kepler:moon}
De Maan, op $r = \qty{3.84e8}{m}$, draait rond in $T = \qty{27.3}{d} =
\qty{2.36e6}{s}$, dus $v = 2\pi r/T \approx \qty{1.02e3}{m/s}$ en
$a = v^2/r \approx \qty{2.7e-3}{m/s^2}$ --- precies $g/60^2$, en de Maan
staat $60$ aardstralen ver: de aantrekking die de appel laat vallen buigt de
Maan, verdund door het omgekeerde kwadraat. Elke seconde ``valt'' de Maan
ongeveer \qty{1.4}{mm} naar ons toe, terwijl haar kilometer opzij haar
eromheen draagt --- de controle die Newton in 1666 uitvoerde.
\end{example}

\section{De wetten van Kepler}

Newton heeft het omgekeerde kwadraat niet geraden: hij haalde het uit drie
regelmatigheden die Johannes Kepler tussen 1609 en 1619 destilleerde uit
twintig jaar planeetposities die Tycho Brahe met het blote oog had gemeten
--- de beste gegevens van de wereld vóór de telescoop.

\begin{definition}[Ellips]\label{def:g12:satellites-kepler:ellipse}
Een \emph{ellips}\index{ellips} is de verzameling punten waarvan de afstanden
tot twee vaste punten, de \emph{brandpunten}\index{brandpunt}, een constante
som hebben; de helft van haar langste middellijn is de \emph{halve lange
as}\index{halve lange as} $a$, en een cirkel is het geval waarin de
brandpunten samenvallen. De punten van de \omterm{def:g10:universal-gravitation:satellite}{omloopbaan} van een planeet die het
dichtst bij en het verst van de Zon liggen, zijn haar
\emph{perihelium}\index{perihelium} en \emph{aphelium}\index{aphelium}.
\end{definition}

\begin{theorem}[Wetten van Kepler]\label{thm:g12:satellites-kepler:kepler}
\index{wetten van Kepler}%
Voor de planeten die om de Zon (massa $M$) draaien:
\begin{enumerate}
  \item elke \omterm{def:g10:universal-gravitation:satellite}{omloopbaan} is een \omterm{def:g12:satellites-kepler:ellipse}{ellips} met de Zon in één \omterm{def:g12:satellites-kepler:ellipse}{brandpunt};
  \item het lijnstuk Zon--planeet bestrijkt gelijke oppervlakken in gelijke
        tijden;
  \item de verhouding van het kwadraat van de \omterm{def:g12:kinematics-2d:ucm}{omlooptijd} tot de derde macht
        van de \omterm{def:g12:satellites-kepler:ellipse}{halve lange as} is voor alle planeten dezelfde:
        $T^2/a^3 = 4\pi^2/(G M)$.
\end{enumerate}
Dezelfde wetten regeren elke familie \omterm{def:g10:universal-gravitation:satellite}{satellieten} om elk centraal lichaam,
met $M$ de centrale massa. \admitted
\end{theorem}

\begin{remark}[Wat we dit jaar kunnen bewijzen]\label{rem:g12:satellites-kepler:proof}
Voor een \emph{cirkelvormige} \omterm{def:g10:universal-gravitation:satellite}{omloopbaan} is de derde wet van ons: kwadrateer
\cref{prop:g12:satellites-kepler:period} en je krijgt
$T^2/r^3 = 4\pi^2/(GM)$ --- één constante voor elke \omterm{def:g10:universal-gravitation:satellite}{satelliet} van hetzelfde
middelpunt. Dat ellipsen, gelijke oppervlakken en de algemene $a$ uit het
omgekeerde kwadraat volgen, wordt met meer analyse in het volume van
bachelorjaar~1 aangetoond. Ondertussen laat de tweede wet zich lezen: de
planeet gaat het snelst in het \omterm{def:g12:satellites-kepler:ellipse}{perihelium}, waar het korte lijnstuk zich moet
haasten.
\end{remark}

\begin{omfigure}
\begin{tikzpicture}[scale=0.85, font=\small,
    declare function={re(\t)=2.025/(1+0.5*cos(\t));}]
  \draw[thick] plot[domain=0:360, samples=120, smooth]
    ({1.35+re(\x)*cos(\x)}, {re(\x)*sin(\x)});
  \fill[omDef!30] (1.35,0) -- plot[domain=-40:40, samples=25]
    ({1.35+re(\x)*cos(\x)}, {re(\x)*sin(\x)}) -- cycle;
  \fill[omDef!30] (1.35,0) -- plot[domain=175:185, samples=10]
    ({1.35+re(\x)*cos(\x)}, {re(\x)*sin(\x)}) -- cycle;
  \fill[omThm] (1.35,0) circle (2.6pt) node[below=2pt] {Zon};
  \draw[-Stealth, omDef, thick]
    ({1.35+re(40)*cos(40)}, {re(40)*sin(40)}) -- ++(-0.55,0.42);
  \draw[-Stealth, omDef, thick]
    ({1.35+re(175)*cos(175)}, {re(175)*sin(175)}) -- ++(-0.12,-0.5);
  \node[right] at (2.75,0) {$P$}; \node[left] at (-2.72,0) {$A$};
\end{tikzpicture}

{\small De tweede wet van Kepler: in één maand bij het \omterm{def:g12:satellites-kepler:ellipse}{perihelium} $P$
bestrijkt de planeet een korte, dikke sector; bij het \omterm{def:g12:satellites-kepler:ellipse}{aphelium} $A$ een lange,
\omterm{def:g11:lenses-and-eye:lens}{dunne} --- gelijke oppervlakken: ze haast zich als ze dichtbij is en treuzelt
als ze ver weg is.}
\end{omfigure}

\begin{omfigure}
\begin{tikzpicture}
\begin{axis}[omaxis, width=8.8cm, height=5.8cm,
    xmode=log, ymode=log,
    xlabel={halve lange as $a$ (astronomische eenheden)},
    ylabel={omlooptijd $T$ (jaren)},
    xmin=0.25, xmax=50, ymin=0.15, ymax=400]
  \addplot[omThm, thick] coordinates {(0.3,0.1643) (40,253)};
  \addplot[only marks, mark=*, mark size=1.7pt, omDef] coordinates
    {(0.387,0.241) (0.723,0.615) (1,1) (1.524,1.881) (5.203,11.86)
     (9.537,29.46) (19.19,84.02) (30.07,164.8)};
  \node[below right, font=\scriptsize] at (axis cs:0.387,0.241) {Mercurius};
  \node[below right, font=\scriptsize] at (axis cs:1,1) {Aarde};
  \node[below right, font=\scriptsize] at (axis cs:5.203,11.86) {Jupiter};
  \node[below right, font=\scriptsize] at (axis cs:30.07,164.8) {Neptunus};
\end{axis}
\end{tikzpicture}

{\small De acht planeten op log--log-assen: een rechte met helling $3/2$,
oftewel $T^2 \propto a^3$ --- de derde wet van Kepler, die in deze eenheden
(astronomische eenheden, jaren) eenvoudig $T^2 = a^3$ luidt.}
\end{omfigure}

\section{Omloopbanen aan het werk}

\begin{definition}[Geostationaire baan]\label{def:g12:satellites-kepler:geo}
Een \omterm{def:g10:universal-gravitation:satellite}{satelliet} beschrijft een
\emph{geostationaire baan}\index{geostationaire baan} wanneer hij boven één
vast punt van de grond blijft hangen. Zijn
\omterm{def:g10:universal-gravitation:satellite}{omloopbaan} moet cirkelvormig zijn, in het vlak van de evenaar liggen, en een
\omterm{def:g12:kinematics-2d:ucm}{omlooptijd} van één \emph{siderische dag}\index{siderische dag} hebben,
$T = \qty{86164}{s}$ (\qty{23}{h} \qty{56}{min} \qty{4}{s}): de tijd die de
Aarde nodig heeft om één keer om haar as te draaien \emph{ten opzichte van
de sterren}. (De dag van \qty{24}{h} is ten opzichte van de Zon, waarheen de
Aarde per dag ook nog een graad opschuift.)
\end{definition}

\begin{example}[De geostationaire hoogte]\label{ex:g12:satellites-kepler:geoalt}
De derde wet van Kepler, achterstevoren gelezen, legt de straal vast:
$r = \left(GMT^2/4\pi^2\right)^{1/3} =
\left(\num{3.98e14} \times (\num{86164})^2/4\pi^2\right)^{1/3}
\approx \qty{4.22e7}{m}$ --- een hoogte $h = r - R \approx
\qty{3.58e7}{m} \approx \qty{35800}{km}$, doorlopen met $v = 2\pi r/T
\approx \qty{3.1}{km/s}$. Elke relais die ``stil hangt'' zit op deze ene
cirkel boven de evenaar --- een andere is er niet.
\end{example}

\begin{example}[De navigatiezwerm]\label{ex:g12:satellites-kepler:gps}
Satellietnavigatiesystemen laten een dertigtal \omterm{def:g10:universal-gravitation:satellite}{satellieten} vliegen op
$r \approx \qty{2.66e7}{m}$ (hoogte $\approx \qty{20200}{km}$), waar
\cref{prop:g12:satellites-kepler:period} $T \approx \qty{4.32e4}{s}$ geeft:
een halve \omterm{def:g12:satellites-kepler:geo}{siderische dag}, zodat elke \omterm{def:g10:universal-gravitation:satellite}{satelliet} zijn spoor over de grond
dagelijks overdoet. Je ontvanger vergelijkt de aankomsttijden van hun
kloksignalen met omloopbanen die tot op \omterm{def:g10:orders-of-magnitude:unit}{meters} bekend zijn.
\end{example}

\begin{method}[Een wereld wegen aan zijn satelliet]\label{met:g12:satellites-kepler:weigh}
\index{een planeet wegen}%
Om de massa van een planeet of een ster te meten: zoek er een willekeurige
\omterm{def:g10:universal-gravitation:satellite}{satelliet} bij (een maan, een sonde, een planeet), meet de straal $r$ en de
\omterm{def:g12:kinematics-2d:ucm}{omlooptijd} $T$ van zijn \omterm{def:g10:universal-gravitation:satellite}{omloopbaan}, en keer de derde wet om:
\[ M = \frac{4\pi^2 r^3}{G\,T^2}. \]
Dit weegt alleen het \emph{centrale} lichaam --- de massa van de \omterm{def:g10:universal-gravitation:satellite}{satelliet}
viel weg in \cref{thm:g12:satellites-kepler:speed}. Elke massa op de
gegevenskaart is zo verkregen: de Aarde uit de Maan, de Zon uit de
\omterm{def:g10:universal-gravitation:satellite}{omloopbaan} van de Aarde zelf.
\end{method}

\begin{remark}[Hoger is trager --- maar duurder]\label{rem:g12:satellites-kepler:energy}
Omdat $v = \sqrt{GM/r}$, kruipt de Maan met \qty{1}{km/s} terwijl het station
met \qty{7.7}{km/s} raast. Toch is de hogere \omterm{def:g10:universal-gravitation:satellite}{omloopbaan} de \emph{duurdere}:
klimmen tegen de \omterm{def:g11:fundamental-interactions:four}{gravitatie} in slaat \omterm{def:g11:mechanical-energy:potential}{potentiële energie} op (de rekening
wordt in \cref{ch:g12:work-and-energy} vereffend), en de winst aan $E_p$
weegt zwaarder dan het verlies aan $E_k$. Daarom brandt een raket
\emph{vooruit} om een hogere, tragere \omterm{def:g10:universal-gravitation:satellite}{omloopbaan} te bereiken --- en spiraalt
een \omterm{def:g10:universal-gravitation:satellite}{satelliet} die door ijle lucht wordt afgeremd \emph{omlaag}, en gaat
daarbij sneller.
\end{remark}

\begin{remark}[Gegevenskaart]\label{rem:g12:satellites-kepler:data}
Tenzij anders vermeld: $G = \qty{6.67e-11}{N.m^2/kg^2}$; Aarde:
$M = \qty{5.97e24}{kg}$, $R = \qty{6.37e6}{m}$, \omterm{def:g12:satellites-kepler:geo}{siderische dag}
\qty{86164}{s}; \omterm{def:g10:universal-gravitation:satellite}{omloopbaan} van de Maan: $r = \qty{3.84e8}{m}$,
$T = \qty{27.3}{d}$; Zon: $M = \qty{1.99e30}{kg}$; afstand Aarde--Zon
\qty{1.496e11}{m}; één jaar $= \qty{3.156e7}{s}$; Mars:
$M = \qty{6.42e23}{kg}$, $R = \qty{3.39e6}{m}$.
\end{remark}

\section{Oefeningen}

\begin{exercise}[$\star$]\label{exo:g12:satellites-kepler:1}
Een \omterm{def:g10:universal-gravitation:satellite}{satelliet} van \qty{1000}{kg} staat op het lanceerplatform en vliegt
daarna op hoogte \qty{400}{km}. Bereken de aantrekking van de Aarde op beide
plaatsen: welk percentage blijft er in de ruimte over --- en waarom zweven de
inzittenden dan?
\end{exercise}

\begin{exercise}[$\star$]\label{exo:g12:satellites-kepler:2}
Bereken voor een (theoretische) \omterm{def:g10:universal-gravitation:satellite}{satelliet} die rakelings over het aardoppervlak
scheert ($r = R$) de baansnelheid en de \omterm{def:g12:kinematics-2d:ucm}{omlooptijd}. Vergelijk met de
kanonskogel van Newton (\cref{rem:g12:satellites-kepler:cannon}).
\end{exercise}

\begin{exercise}[$\star$]\label{exo:g12:satellites-kepler:3}
Een \omterm{def:g10:universal-gravitation:satellite}{satelliet} wordt van een cirkelvormige \omterm{def:g10:universal-gravitation:satellite}{omloopbaan} met straal $r$ naar een
met straal $2r$ gebracht. Met welke factor veranderen zijn snelheid en zijn
\omterm{def:g12:kinematics-2d:ucm}{omlooptijd}? Hangen de antwoorden van zijn massa af?
\end{exercise}

\begin{exercise}[$\star$]\label{exo:g12:satellites-kepler:4}
Doe de getallen van het station over: bereken uit $r = \qty{6.77e6}{m}$ de
waarden van $v$ en $T$, en het aantal omloopbanen dat het in \qty{24}{h}
aflegt.
\end{exercise}

\begin{exercise}[$\star$]\label{exo:g12:satellites-kepler:5}
Formuleer de drie \omterm{thm:g12:satellites-kepler:kepler}{wetten van Kepler}. Waar op zijn zeer uitgerekte \omterm{def:g12:satellites-kepler:ellipse}{ellips} gaat
een komeet het snelst, en welke wet zegt dat? Hangt de \omterm{def:g12:kinematics-2d:ucm}{omlooptijd} van een
\omterm{def:g10:universal-gravitation:satellite}{satelliet} van zijn eigen massa af?
\end{exercise}

\begin{exercise}[$\star\star$]\label{exo:g12:satellites-kepler:6}
Weeg de Aarde: bereken uit de straal en de \omterm{def:g12:kinematics-2d:ucm}{omlooptijd} van de Maan
(gegevenskaart) de massa van de Aarde en vergelijk met de waarde op de kaart.
\end{exercise}

\begin{exercise}[$\star\star$]\label{exo:g12:satellites-kepler:7}
Een bedrijf stelt een \omterm{def:g10:universal-gravitation:satellite}{satelliet} voor die permanent boven Parijs
(breedtegraad $49^\circ$ noord) blijft hangen. Leg met de richting van de
\omterm{def:g11:fundamental-interactions:four}{gravitatie} uit waarom elke cirkelvormige \omterm{def:g10:universal-gravitation:satellite}{omloopbaan} het middelpunt van de
Aarde als middelpunt heeft; concludeer dat het voorstel onmogelijk is. Waar
\emph{kan} een \omterm{def:g10:universal-gravitation:satellite}{satelliet} blijven hangen?
\end{exercise}

\begin{exercise}[$\star\star$]\label{exo:g12:satellites-kepler:8}
Phobos draait om Mars op $r = \qty{9.38e6}{m}$ in $T = \qty{7}{h}
\qty{39}{min}$. Leid de massa van Mars af; vergelijk met de gegevenskaart.
\end{exercise}

\begin{exercise}[$\star\star$]\label{exo:g12:satellites-kepler:9}
Mars draait om de Zon op $a = \qty{2.279e11}{m}$ in \qty{687}{d}. Bereken
$T^2/a^3$ voor Mars en voor de Aarde (gegevenskaart), controleer de derde wet
van Kepler, en leid de massa van de Zon af.
\end{exercise}

\begin{exercise}[$\star\star$]\label{exo:g12:satellites-kepler:10}
De maantoets van Newton: bereken uit de gegevenskaart de snelheid van de Maan
en haar \omterm{thm:g12:kinematics-2d:centripetal}{middelpuntzoekende versnelling} $a = v^2/r$. Toon aan dat die gelijk
is aan $g/60^2$, en zeg waarom dat Newton ervan overtuigde dat één wet over
appel en Maan regeert.
\end{exercise}

\begin{exercise}[$\star\star$]\label{exo:g12:satellites-kepler:11}
In eenheden van de Zon (astronomische eenheden, jaren) luidt de derde wet van
Kepler $T^2 = a^3$: de rechte van de log--log-grafiek uit de cursus. Zet er de
planetoïde Ceres op, $a = \qty{2.77}{au}$, en de komeet van Halley,
$a = \qty{17.8}{au}$: bereken beide \omterm{def:g12:kinematics-2d:ucm}{omlooptijden}. Halley kwam voor het laatst
in 1986 door het \omterm{def:g12:satellites-kepler:ellipse}{perihelium} --- wanneer wordt hij terugverwacht?
\end{exercise}

\begin{exercise}[$\star\star\star$]\label{exo:g12:satellites-kepler:12}
IJle lucht strijkt langs het station, en toch \emph{neemt} zijn snelheid
\emph{toe} naarmate het hoogte verliest. Bereken $v$ op de hoogten
\qty{400}{km} en \qty{350}{km}, en los daarna de paradox op: \omterm{def:g10:inertia:inventory}{wrijving} remt
dingen af --- waar komt de extra snelheid vandaan?
\end{exercise}

\begin{exercise}[$\star\star\star$]\label{exo:g12:satellites-kepler:13}
Een telescoop vindt een exoplaneet die elke \qty{3.50}{d} om haar ster draait
op $r = \qty{7.48e9}{m}$. Weeg de ster, in \omterm{def:g10:orders-of-magnitude:unit}{kilogram} en in zonsmassa's. Welke
massa --- die van de ster of die van de planeet --- valt zo \emph{niet} te
krijgen, en waarom?
\end{exercise}

\begin{exercise}[$\star\star\star$]\label{exo:g12:satellites-kepler:14}
Kolonisten op Mars willen een ``areostationair'' relais boven hun basis. Mars
draait ten opzichte van de sterren in \qty{24}{h} \qty{37}{min} om haar as:
bereken de straal van de \omterm{def:g10:universal-gravitation:satellite}{omloopbaan} en de hoogte boven het marsoppervlak
(gegevenskaart).
\end{exercise}

\begin{exercise}[$\star\star\star$]\label{exo:g12:satellites-kepler:15}
De komeet van Halley passeert het \omterm{def:g12:satellites-kepler:ellipse}{perihelium} op $r_p = \qty{8.77e10}{m}$ met
\qty{54.5}{km/s}; het \omterm{def:g12:satellites-kepler:ellipse}{aphelium} ligt op $r_a = \qty{5.25e12}{m}$. In beide
uitersten staat de snelheid loodrecht op het lijnstuk Zon--komeet, en de
tweede wet van Kepler dwingt dan $r_p v_p = r_a v_a$ af (gelijke bestreken
driehoeken met oppervlakte $\tfrac12 r v\,\Delta t$). Leid de snelheid in het
\omterm{def:g12:satellites-kepler:ellipse}{aphelium} af en becommentarieer de verhouding.
\end{exercise}

\section{Probleem: Een plaats boven de evenaar}

\begin{problem}[{Weekendopgave --- een \omterm{def:g10:universal-gravitation:satellite}{satelliet} boven de evenaar zetten:
waarom de enige parkeerplaats aan de hemel een cirkel op \qty{35800}{km}
hoogte is, hoe die zich verhoudt tot het station dat eronder voorbijraast, en
hoe dezelfde wet onderweg Jupiter weegt}]
\label{pb:g12:satellites-kepler:1}
Een telecombedrijf wil een relais waar zijn schotels, eenmaal vastgeschroefd,
nooit achteraan hoeven te jagen: een geostationaire \omterm{def:g10:universal-gravitation:satellite}{satelliet}. Jij bent de
missieanalist; gebruik de gegevenskaart. Voor deel~IV: de maan Io van Jupiter
draait op $r = \qty{4.22e8}{m}$ in $T = \qty{1.77}{d}$, en Jupiter draait ten
opzichte van de sterren in \qty{9}{h} \qty{56}{min} om haar as.

\medskip
\noindent\textbf{Deel I --- De regels van het stilhangen.}
\begin{enumerate}
  \item De \omterm{def:g10:universal-gravitation:satellite}{satelliet} moet boven één vast punt van de grond blijven. Wat moet
        zijn \omterm{def:g12:kinematics-2d:ucm}{omlooptijd} zijn --- en waarom de \omterm{def:g12:satellites-kepler:geo}{siderische dag} van
        \qty{86164}{s} en niet \qty{24}{h}?
  \item Bij een cirkelbeweging wijst de \omterm{def:g11:forces-and-motion:netforce}{resulterende kracht} naar het
        middelpunt van de cirkel; hier is de \omterm{def:g11:fundamental-interactions:four}{gravitatie} de enige \omterm{def:g10:inertia:force}{kracht}. Leid
        af dat het middelpunt van elke cirkelvormige \omterm{def:g10:universal-gravitation:satellite}{omloopbaan} het
        middelpunt van de Aarde is.
  \item Een cirkel die boven de breedtegraad van Parijs hangt, zou zijn
        middelpunt op de as van de Aarde hebben, ten noorden van het
        middelpunt. Concludeer: in welk vlak moet een \omterm{def:g12:satellites-kepler:geo}{geostationaire baan}
        liggen?
  \item Schrijf de tweede wet van Newton in het \omterm{def:g12:kinematics-2d:frenet}{frenetstelsel} op voor een
        cirkelvormige \omterm{def:g10:universal-gravitation:satellite}{omloopbaan} met straal $r$ en leid $v = \sqrt{GM/r}$ af.
  \item Leid $T = 2\pi\sqrt{r^3/(GM)}$ af en herschrijf die als de derde wet
        van Kepler, $T^2/r^3 = 4\pi^2/(GM)$.
\end{enumerate}

\medskip
\noindent\textbf{Deel II --- De ene cirkel die werkt.}
\begin{enumerate}[resume]
  \item Los de derde wet op naar $r$ en bereken de geostationaire straal.
  \item Leid daaruit de hoogte boven de grond af.
  \item Bereken de baansnelheid uit $v = 2\pi r/T$.
  \item Controleer die met $v = \sqrt{GM/r}$.
  \item De \omterm{def:g10:universal-gravitation:satellite}{satelliet} van het bedrijf heeft massa \qty{4.0e3}{kg}; die van een
        concurrent weegt het dubbele. Vergelijk hun geostationaire stralen, en
        zeg wat de extra massa bij de lancering verandert.
\end{enumerate}

\medskip
\noindent\textbf{Deel III --- Het station eronder.}
\begin{enumerate}[resume]
  \item Het ruimtestation vliegt op $r = \qty{6.77e6}{m}$: bereken zijn
        snelheid.
  \item Bereken zijn \omterm{def:g12:kinematics-2d:ucm}{omlooptijd}, in seconden en in minuten.
  \item Hoeveel omloopbanen legt het in \qty{24}{h} af --- ongeveer hoeveel
        zonsopgangen zien de astronauten per dag?
  \item Bereken de verhouding van de stralen
        $r_{\text{geo}}/r_{\text{ISS}}$, en ga na dat de verhouding van de
        \omterm{def:g12:kinematics-2d:ucm}{omlooptijden} dat getal tot de macht $3/2$ is, zoals de derde wet van
        Kepler eist.
  \item Welke \omterm{def:g10:universal-gravitation:satellite}{satelliet} gaat sneller, en welke kostte meer \omterm{def:g10:energy-conservation:energy}{energie} per
        \omterm{def:g10:orders-of-magnitude:unit}{kilogram} om op zijn plaats te krijgen? Verzoen de twee in één zin.
\end{enumerate}

\medskip
\noindent\textbf{Deel IV --- Jupiter er en passant bij wegen.}
\begin{enumerate}[resume]
  \item \omterm{def:g10:pressure:pressure}{Druk} met de derde wet van Kepler, toegepast op Io, de massa $M_J$ van
        Jupiter uit in $G$, $r$ en $T$.
  \item Bereken $M_J$.
  \item Hoeveel aardmassa's is dat? Ruwweg welke fractie van de massa van de
        Zon?
  \item Een ``joviestationair'' relais zou boven één punt van de wolken van
        Jupiter moeten hangen: bereken de straal van zijn \omterm{def:g10:universal-gravitation:satellite}{omloopbaan};
        vergelijk met die van Io.
  \item Rapporteer aan de directie, in twee zinnen: de hoogte waar het bedrijf
        zijn relais moet parkeren, en de massa van Jupiter die dezelfde wet er
        gratis bij heeft geleverd.
\end{enumerate}
\end{problem}
